%% ====================================================================
%%  SECTION 1 --- INTRODUCTION (~2.5pp)
%% ====================================================================
\section{Introduction}
\label{sec:intro}

A builder of agent systems faces a question every deployment cycle: \emph{I have a task family, a compute budget, and a menu of models ranging from fine-tuned specialists to frontier generalists.
What system should I build?}

The question is not merely which model to call.
It is which model, which orchestration policy, and how many sequential attempts jointly minimize time to verified solution.
The dominant methodology is benchmark racing: run both models on a test set, pick the winner, repeat when the next model arrives.

Benchmark racing is fast and requires no theory.
It is also structurally blind.
It conflates per-step cost, within-mode competence, task-relevant information, and routing quality into a single opaque comparison.
When GPT-4 beats Claude on one benchmark but Claude beats GPT-4 on another, benchmark racing cannot diagnose whether the difference comes from per-step competence, from routing mismatch (the evaluation protocol suits one model's strengths), from cost allocation, or from the inherent information structure of the task.
All four factors are real; they interact, and benchmark racing conflates them into a single number.

This paper disentangles them.

\paragraph{The decomposition.}
Under the packed latent-mode family assumption (a task distribution with discrete solution strategies), the performance gap between any two agent designs decomposes exactly as
\begin{equation}
\label{eq:four-term-intro}
\Delta = \log\frac{\kappa_0}{\kappa} + \varphi + G - \varepsilon,
\end{equation}
separating per-step cost, within-mode competence, information generation quality, and routing mismatch into four independently measurable terms.
The decomposition is an algebraic identity---no approximation, no limit.
A builder who estimates these four numbers from a structured pilot knows which lever to pull: switch models if cost or competence dominates, rebuild the router if mismatch~$\varepsilon$ is large, or restructure the task distribution if the information term~$G$ is small.

The hero equation is simpler still.
With model held fixed, the Routing-Information Identity reduces to
\[
\Delta = G - \varepsilon:
\]
certified log-speedup equals information generated minus information wasted.

A natural expectation, given the algebraic structure, is that these four terms are shadow prices of resource constraints---the dual variables of a Lagrangian relaxation.
We show that this expectation is precisely wrong (Section~\ref{sec:lagrangian}).
The four terms are components of the primal objective, not the dual.
The decomposition's value to optimization is as a separable evaluator, not as a dual decomposition.

\paragraph{The crossover theorem.}
The decomposition also yields a concrete decision rule.
Given a large expensive model (success probability $p_l$ per attempt, cost $\kappa_l$) and a small cheap model ($p_s$, $\kappa_s$), the number of cheap-model attempts needed to match one expensive-model attempt is
\begin{equation}
\label{eq:crossover-intro}
\Dcross = \frac{\log(1 - p_l)}{\log(1 - p_s)}.
\end{equation}
If $\Dcross \le \kappa_l/\kappa_s$, the small model wins.
This is the paper's most actionable result: one formula, immediately deployable.

\paragraph{Challenging folklore.}
The paper directly challenges four widespread assumptions:

\begin{enumerate}
\item \emph{``Scale is all you need.''}\;
Model competence is one of four terms, and not always the dominant one.

\item \emph{``More agents is better.''}\;
Optimal orchestration depth \emph{decreases} with model competence: for strong models, the optimal architecture is simpler.

\item \emph{``Benchmark performance predicts deployment performance.''}\;
The decomposition shows exactly what benchmarks miss: routing mismatch and information structure are invisible to fixed-routing evaluation.

\item \emph{``Agent design is engineering, not science.''}\;
Agent design has mathematical structure amenable to theory---simpler than optimization theory, but real.
\end{enumerate}

\paragraph{Contributions.}
The paper makes five contributions:
(1)~An exact four-term decomposition of the performance gap between agent designs (Section~\ref{sec:decomposition}), valid under the packed-family assumption with a bounded remainder outside it.
(2)~An honest negative result: the four terms are performance accounting, not Lagrangian shadow prices (Section~\ref{sec:lagrangian}).
(3)~A graph-depth model characterizing optimal orchestration depth, with monotonicity in competence and cost (Section~\ref{sec:graph-depth}), and a closed-form crossover theorem (Section~\ref{sec:crossover}).
(4)~Proxy validity (Section~\ref{sec:proxy-validity}) and Bernstein-improved sample complexity bounds (Section~\ref{sec:sample-complexity}) that make the framework practical.
(5)~Synthetic validation showing up to 50x evaluation efficiency over benchmark racing in settings with $k = 100$ designs (Section~\ref{sec:validation}), with honest scoping of when the framework should not be used (Section~\ref{sec:scoping}).

\paragraph{Roadmap.}
Section~\ref{sec:prelim} defines the mathematical objects.
Sections~\ref{sec:decomposition}--\ref{sec:lagrangian} develop the decomposition and its Lagrangian non-connection.
Sections~\ref{sec:graph-depth}--\ref{sec:bernstein} derive the depth-cost tradeoff and sample complexity.
Section~\ref{sec:validation} validates and scopes the framework.
Section~\ref{sec:related} positions the paper.
Section~\ref{sec:discussion} discusses limitations and open problems.
